\documentclass[../../../include/open-logic-section]{subfiles}

\begin{document}

\iftag{FOL}
      {\olfileid{fol}{axd}{ded}}
      {\olfileid{pl}{axd}{ded}}

\olsection{演绎定理}

% Part: first-order-logic
令 $\Gamma$ 为一个语句集，并令 $!A$ 和 $!B$ 为语句。
则 $\Gamma \cup \{!A\} \Proves !B$ 当且仅当 $\Gamma \Proves !A \lif !B$。
% Chapter: axiomatic-deduction

% Section: deduction-theorem
概略地说，这个定理表达的意思是：要由假设 $!A$ 证明 $!B$，只需证明那个条件句 $!A \lif !B$ 成立。
% verification of properties of provability needed for maximally
正是这种能够从“由 $!A$ 证明 $!B$”转到“证明 $!A \lif !B$”并能转回来的能力，使得条件证明法成为一种可行的证明方法。
% consistent sets in the completeness chapter.
它是该推导系统的一项元理论性质，为条件证明法的使用提供了依据。

我们已经看到，在公理系统中进行推导是繁琐的，且往往难以找到推导。相较于写出完整的一长串公式，通常更简单的做法是利用几个简单的结论来论证推导的存在。我们已经得到了三个这样的结论：\olref[ptn]{prop:reflexivity} 表明：若已知 $!A \in \Gamma$，则可断言 $\Gamma \Proves !A$。\olref[ptn]{prop:monotonicity} 表明：若 $\Gamma \Proves !A$，则也有 $\Gamma \cup \{!B\} \Proves !A$。而 \olref[ptn]{prop:transitivity} 表明：若 $\Gamma \Proves !A$ 且 $!A \Proves !B$，则 $\Gamma \Proves !B$。下面是另一个简单的结论，即肯定前件规则的“元”版本：

\begin{prop}
  \ollabel{prop:mp} 若 $\Gamma \Proves !A$ 且 $\Gamma \Proves !A \lif !B$，则 $\Gamma \Proves !B$。
\end{prop}

\begin{proof}
  我们有 $\{!A, !A \lif !B\} \Proves !B$：
  \begin{derivation}
    1. & $!A$ & 假设\\
    2. & $!A \lif !B$ & 假设\\
    3. & $!B$ & 1, 2, MP
  \end{derivation}
  根据 \olref[ptn]{prop:transitivity}，可得 $\Gamma \Proves !B$。
\end{proof}

在此语境下我们将使用的最重要的结果是演绎定理：

\begin{thm}[演绎定理]
\ollabel{thm:deduction-thm} $\Gamma \cup \{!A\} \Proves !B$ 当且仅当 $\Gamma \Proves !A \lif !B$。
\end{thm}

\begin{proof}
“如果”方向是显然的。如果 $\Gamma \Proves !A \lif !B$，那么由
\olref[ptn]{prop:monotonicity} 可得 $\Gamma \cup \{!A\} \Proves !A \lif !B$。同时，由
\olref[ptn]{prop:reflexivity} 可得 $\Gamma \cup \{!A\} \Proves !A$。因此，根据 \olref{prop:mp}，有 $\Gamma \cup
\{!A\} \Proves !B$。

对于“仅当”方向，我们对从 $\Gamma \cup \{!A\}$ 到 $!B$ 的推导长度进行归纳。

对于归纳基础，我们证明结论对每个长度为~$1$ 的推导成立。一个从 $\Gamma \cup \{!A\}$ 到 $!B$、长度为~$1$ 的推导只包含 $!B$ 本身；且既然该推导是合法的，$!B$ 要么属于 $\Gamma \cup \{!A\}$，要么是一条公理。如果 $!B \in \Gamma$ 或是公理，那么 $\Gamma \Proves !B$。又由 \olref[prp]{ax:lif1} 可得 $\Gamma
\Proves !B \lif (!A \lif !B)$，于是根据 \olref{prop:mp} 有 $\Gamma \Proves !A \lif !B$。如果 $!B \in \{ !A\}$，那么 $\Gamma \Proves !A \lif !B$，因为此时 $!A \lif !B$ 就是 $!A \lif !A$，而 \olref[pro]{ex:identity} 中已给出了该公式的推导。

对于归纳步骤，假设从 $\Gamma \cup \{!A\}$ 到 $!B$ 的推导，其最后一步 $!B$ 是由肯定前件规则得出的。（如果最后一步不是由肯定前件规则得出，则 $!B \in \Gamma$，$!B \ident !A$，或者 $!B$ 是公理，此时可应用与归纳基础相同的推理。）那么，在该推导中存在某个公式 $!C$，使得前面的步骤包含 $!C \lif !B$ 和 $!C$，即 $\Gamma \cup \{!A\} \Proves !C \lif !B$ 和 $\Gamma \cup \{!A\}
\Proves !C$，并且它们各自的推导长度更短，因此可以对它们应用归纳假设。于是我们得到：
\begin{align*}
  & \Gamma \Proves !A \lif (!C \lif !B); \\
  & \Gamma \Proves !A \lif !C.
\end{align*}
此外，由 \olref[prp]{ax:lif2} 可得
\[
\Gamma \Proves (!A \lif (!C \lif !B)) \lif
((!A\lif !C)  \lif (!A \lif !B)),
\]
再应用两次 \olref{prop:mp} 即得 $\Gamma \Proves !A \lif !B$，如所求。
\end{proof}

请注意，\olref[prp]{ax:lif1} 和 \olref[prp]{ax:lif2} 的选取正是为了保证演绎定理成立。

以下是一些关于可推导性的有用结论，我们将其留作练习。

\begin{prop}
  \ollabel{prop:derivfacts}
\begin{enumerate}
\item $\Proves (!A \lif !B) \lif ((!B \lif !C)
  \lif (!A \lif !C)$; \ollabel{derivfacts:a}
\item 若 $\Gamma \cup \{ \lnot !A\}
  \Proves \lnot !B$ 则 $\Gamma \cup \{ !B\} \Proves
  !A$（假言易位）; \ollabel{derivfacts:b}
\item  $\{ !A, \lnot!A\} \Proves
    !B$（爆炸律）; \ollabel{derivfacts:c}
\item  $\{ \lnot\lnot!A\} \Proves
  !A$（双重否定消除）;\ollabel{derivfacts:d}
\item 若 $\Gamma \Proves \lnot\lnot!A$ 则 $\Gamma \Proves
  !A$;\ollabel{derivfacts:e}
\end{enumerate}
\end{prop}

\tagprob{FOL}
\begin{prob}
  证明 \olref[fol][axd][ded]{prop:derivfacts}
\end{prob}
\tagendprob

\tagprob{notFOL}
\begin{prob}
  证明 \olref[pl][axd][ded]{prop:derivfacts}
\end{prob}
\tagendprob

\end{document}